% !TeX root = ../../Book3.tex
\section{双曲测地线}
\label{b3:sec:1302}

圆盘的最短路是直径的自同构像。完备性来自边界的无限距离，而常负曲率则描述同一局部尺度的微分几何。

\subsection{测地线}
\label{b3:subsec:130201}

\begin{theorem}\label{b3:thm:13-disk-geodesics}
圆盘中任意不同两点由唯一的测地线段连接，忽略保持方向的重参数化。测地线是与单位圆正交的圆弧或直径。
\end{theorem}
\begin{proof}
以自同构把一端移到0。距离证明中的两个不等式同时取等，要求路径速度几乎处处径向且半径单调，故只有径向线段能最短。撤销自同构后得到唯一线段。Möbius 变换保持广义圆与夹角，直径与单位圆正交，所以像为正交圆弧。反向每条这样的圆弧可由一个圆盘自同构送成直径，故确为测地线。
\end{proof}
沿直径可用弧长参数 \(r=\tanh(s/2)\)，\(-\infty<s<\infty\)，说明完整测地线的两端都在圆周，而不会在有限弧长到达边界。

\subsection{完备性}
\label{b3:subsec:130202}

\begin{theorem}\label{b3:thm:13-disk-complete}
圆盘 Poincaré 度量完备，且闭有界距离球为紧集。
\end{theorem}
\begin{proof}
0点距离公式使半径 \(M\) 的闭球为欧氏闭盘 \(|z|\le\tanh(M/2)<1\)，故紧；其他中心由自同构传回。Cauchy 列到某个固定点的距离有界，因此落在这样的紧球。密度至少2，故它也是欧氏 Cauchy 列，并有盘内极限；在包含该球的略大紧盘上密度有上界，直线段估计使它在 Poincaré 距离也趋于这个极限。
\end{proof}
完备不等于底空间欧氏闭合。开圆盘在欧氏空间不完备，在这个长度单位中却完备。度量改变了 Cauchy 列的条件，使逃向边界的点列不再 Cauchy。

\begin{proposition}[双曲圆的长度与面积]\label{b3:prop:13-ball-area}
双曲半径 \(R\) 的圆盘面积与边界长度分别为
\[
 A(R)=2\pi(\cosh R-1)=4\pi\sinh^2(R/2),\qquad
 L(R)=2\pi\sinh R.
\]
\end{proposition}
\begin{proof}
由自同构不变性，只需中心0。双曲半径 \(R\) 对应 Euclid 半径 \(r=\tanh(R/2)\)。共形面积元是 \(\lambda^2\,\mathrm{d}A\)，故
\[
 A=\int_0^r\frac{8\pi t}{(1-t^2)^2}\,\mathrm{d}t
   =\frac{4\pi r^2}{1-r^2}=4\pi\sinh^2(R/2).
\]
圆周上密度恒为 \(2/(1-r^2)\)，长度为 \(4\pi r/(1-r^2)=2\pi\sinh R\)。
\end{proof}
小半径时 \(A(R)\sim\pi R^2\)、\(L(R)\sim2\pi R\)，符合局部 Euclid 行为；大半径时两者均指数增长。完整双曲盘的总面积无穷，而每个有限半径闭球紧。这两项性质相容，完备性不意味着总面积有限。

\ProseParagraphBreak
闭球紧也可从距离公式直接看。以0为中心的闭球是 \(|z|\le\tanh(R/2)<1\) 的 Euclid 闭盘；任意中心由自同构化到0。每条有界距离的序列因此有收敛子列，这比仅有 Cauchy 完备性更强。后面甲板轨道的距离下确界能够达到，正是依赖闭球紧与轨道离散共同作用。

\subsection{常负曲率}
\label{b3:subsec:130203}

对共形长度元 \(\lambda|\mathrm{d}z|\)，Gaussian 曲率计算式为
\[
 K=-\lambda^{-2}\Delta\log\lambda.
\]
令 \(u=\log\lambda\)，度量张量为 \(e^{2u}(\mathrm{d}x^2+\mathrm{d}y^2)\)。其 Christoffel 系数为 \(\Gamma^1_{11}=u_x,\Gamma^1_{12}=u_y,\Gamma^1_{22}=-u_x\)，另三项由对称交换坐标得到；代入曲率张量后，一次导数的平方项相消，留下 \(-e^{-2u}(u_{xx}+u_{yy})\)。因此上述式子没有使用统一化。

\ProseParagraphBreak
对圆盘 \(u=\log2-\log(1-|z|^2)\)，直接计算 \(\Delta u=4/(1-|z|^2)^2=\lambda^2\)，所以 \(K=-1\)。常数密度的欧氏度量曲率0，不能通过整体缩放变成负曲率；位置变化的密度是决定性因素。

\subsection{等距变换}
\label{b3:subsec:130204}

\begin{proposition}\label{b3:prop:13-disk-isometries}
保定向光滑 Poincaré 等距变换恰为圆盘 Möbius 自同构；反定向光滑等距变换是它们与复共轭的复合。
\end{proposition}
\begin{proof}
等距条件使导数在每点为欧氏相似变换。保定向时导数满足 Cauchy--Riemann 方程，故变换全纯；它及逆均全纯，因此由圆盘自同构分类为 Möbius。反定向时先复合共轭转成保定向。反向由密度不变性和共轭不变性得到。
\end{proof}
这里给出光滑 Riemannian 等距变换的分类，没有把一般度量空间等距映射的自动光滑性作为未证明前提。

\ProseParagraphBreak
径向线虽都经过0，通常的 Euclid 直线段却不一定是双曲测地线。两点若不在同一直径上，连接它们的最短曲线是与单位圆正交的圆弧。圆弧可以通过将第一点送到0的 Möbius 变换取得，再把径向段拉回。Möbius 变换保持圆及直线的广义圆，也保持交角，所以得到正交性。弧的唯一性来自拉回后径向最短路的唯一性，而不是仅凭图中能画出一条圆弧。

\ProseParagraphBreak
圆盘自同构可能没有内部不动点，这不妨碍它为等距变换。抛物型和双曲型的不动点在边界，而边界处距离无穷，不属于度量空间。椭圆型有内部不动点，因此不能作为曲面无分歧覆盖的非平凡甲板变换。这把自同构的动力分类与自由覆盖作用联系起来。
